% Abhakseite „Das kann ich“ – Zone nur im Gesamt (\mitzone)
%% TODO Ich-kann-Sätze: Platzhalter wie in den Aufgabendateien
\begin{abhakseite}
\ifdefined\mitzone
\abhakgruppe{Kennst du schon}
\abhak{1}{Kleines Einmaleins und Malnehmen mehrstelliger Zahlen im Kopf oder halbschriftlich, Malketten mit gleichen Faktoren}
\abhak{2}{Quadratzahlen bis 100 aus dem Kopf, Kubikzahlen zwei hoch drei, drei hoch drei, zehn hoch drei}
\abhak{3}{Vorzeichenregel beim Malnehmen (minus mal minus gibt plus), Anzahl der Minuszeichen zählen, Quadrat einer negativen Zahl mit und ohne Klammer}
\abhak{4}{Dezimalzahlen multiplizieren (Kommastellen zählen), Brüche multiplizieren}
\abhak{5}{Dezimalzahlen und Brüche vergleichen und ordnen, Bruch in Dezimalzahl umwandeln, Prozent in Dezimalzahl, negative Zahlen auf der Zahlengerade ordnen}
\abhak{6}{Taschenrechner}
\abhak{7}{Große Zahlen lesen und schreiben (Dreiergruppen, Zahlwörter bis Billion), Stellenwerttafel, mal zehn, hundert und tausend durch Nullen anhängen}
\abhak{8}{Kommaverschiebung}
\abhak{9}{Runden von Dezimalzahlen auf eine und zwei Stellen, Näherungswert mit ≈ schreiben}
\abhak{10}{Vorzeichenregel beim Malnehmen (minus mal minus gibt plus), Anzahl der Minuszeichen zählen, Quadrat einer negativen Zahl mit und ohne Klammer – Fehler finden}
\abhak{11}{Vorzeichenregel beim Malnehmen (minus mal minus gibt plus), Anzahl der Minuszeichen zählen, Quadrat einer negativen Zahl mit und ohne Klammer – selbst rechnen}
\fi
\abhakgruppe{Einheit 1 · Potenzen}
\abhak{12}{Plus oder minus?}
\abhak{13}{Bruch oder minus?}
\abhak{14}{Potenzen}
\abhak{15}{Potenzen von Brüchen}
\abhak{16}{Potenz gegen Dezimalzahl und Prozent vergleichen}
\abhak{17}{hoch null}
\abhak{18}{Potenzen – Fehler finden, Begründen}
\abhakgruppe{Einheit 2 · Zehnerpotenzen}
\abhak{19}{Zehnerpotenzen}
\abhak{20}{Zehnerpotenzen – Fehler finden, Begründen}
\abhakgruppe{Einheit 3 · Quadratwurzeln}
\abhak{21}{Was zuerst?}
\abhak{22}{Quadratwurzeln}
\abhak{23}{Quadratwurzeln – Fehler finden, Begründen}
\end{abhakseite}
